\documentclass[12pt]{article}
\usepackage[utf8]{inputenc}
\usepackage[T1]{fontenc}
\usepackage{hyperref}
\title{Wajarri Language Resources for K-12 Curriculum Development}
\author{Compiled by Claude Sonnet 5 (claude-sonnet-5) for Andr\'as Kornai}
\date{}
\begin{document}
\maketitle

\section{Introduction}
Wajarri (also spelled Wadjarri or Watjarri; ISO 639-3 wbv) is a Pama-Nyungan language of the Murchison region of Western Australia. Searches found a few dictionary and learning resources, mostly from the Geraldton-based Irra Wangga language programme. They also found two unpublished or academic sketch grammars. No school curriculum was found.

\section{Primary Educational Resources}
The Bundiyarra Irra Wangga programme is reported to hold educational material, including books, posters and flashcards. No online catalogue was found, so availability should be confirmed with the programme directly.
\begin{itemize}
\item YMAC news item on the dictionary launch, which also mentions the educational materials: \url{https://www.ymac.org.au/official-launch-of-the-wadjarri-language-dictionary/}
\end{itemize}

\section{Linguistic and Reference Materials}
\begin{itemize}
\item Wadjarri Language Dictionary (2012), produced by the Yamaji Language Centre and Bundiyarra Irra Wangga Language Program. The launch report describes it as taking over 25 years to compile and as a 2500-word dictionary: \url{https://www.ymac.org.au/official-launch-of-the-wadjarri-language-dictionary/}
\item YMAC tag page for the Wajarri dictionary: \url{https://www.ymac.org.au/tag/wajarri-dictionary/}
\item Sketch grammars by Douglas (1981) and Marmion (1996, unpublished Honours thesis, University of New England). These are cited in the Wikipedia article; no online copy was seen: \url{https://en.wikipedia.org/wiki/Wajarri_language}
\end{itemize}

\section{Religious and Community Resources}
\begin{itemize}
\item Yamaji Language Centre (YMAC tag page): \url{https://www.ymac.org.au/tag/yamaji-language-centre/}
\item Irra Wangga (Geraldton) Language Program (YMAC tag page): \url{https://www.ymac.org.au/tag/irra-wangga-geraldton-language-program/}
\end{itemize}

\section{Digital and Audio Resources}
Wikipedia states that the Irra Wangga programme has produced a dictionary app. No app store page or audio archive was seen in the search results: \url{https://en.wikipedia.org/wiki/Wajarri_language}

\section{Curriculum Development Context}
The language is severely endangered, with about 20 speakers. A K-12 curriculum would depend on the Irra Wangga programme and the Yamaji Language Centre. The dictionary is the main documented base, and the grammars are sketch-level. No school programme or published curriculum was found.
\end{document}